%% =========================================================================
%%  Приложение: лемма выравнивания направлений в VIJI-итерациях
%%
%%  Использует обозначения главы 4 (sections/ch4_sp_afd.tex) и теоремные
%%  окружения / макросы (\Norm, \Fro, \Jstar, ...) из main.tex.
%% =========================================================================

\chapter{Выравнивание направлений в VIJI-итерациях}
\label{app:alignment}

В настоящем приложении доказывается лемма~\ref{lem:alignment} для
$\mu$-сильно монотонного $L_1$-гладкого $F$ при единственном
дополнительном условии~--- асимптотическом выравнивании
направлений ошибки \textup{(}\ref{as:dir-align}\textup{)},
см.~\S\ref{sec:app-aln-l1:secant-rk}. Стандартная Q-линейная сходимость
VIJI-итерации \textup{(}\cite[Theorem~3.4]{Agafonov2024VIJI}\textup{)}
используется как фоновая. Все константы зависят только от
$\mu, L_1, \eta$ и равномерной оценки $M_0 := \sup_k \Fro{B_k - \Jstar} <
\infty$ из следствия~\ref{cor:sp-in-segment}.

\paragraph{Обозначение $O(\cdot)$.}
Везде $g = O(h)$ означает $\Norm{g} \le C\,h$ с константой $C \ge 0$,
зависящей только от $\mu, L_1, \eta, M_0, c_d$ \textup{(}при использовании
\ref{as:step-ndg}\textup{)}, $\rho_v$ \textup{(}Q-линейная константа из
\cite[Theorem~3.4]{Agafonov2024VIJI}\textup{)} и $\omega_h$
\textup{(}при использовании \ref{as:dir-align}\textup{)}, но не от номера
итерации $k$ или сегмента~$s$.

\section{Постановка и линеаризация}
\label{sec:app-aln-l1:setup}

\paragraph{Разложение оператора.}
Положим $A := \Jstar = \nabla F(x^{*})$ и
\begin{equation}
F(x) \;=\; A\,(x - x^{*}) \;+\; R(x),
\qquad
R(x^{*}) = 0,
\quad
\nabla R(x^{*}) = 0.
\label{eq:app-aln-l1:decomp}
\end{equation}
По $L_1$-липшицевости $\nabla F$
\begin{equation}
\Norm{R(x) - R(y) - \nabla R(y)\,(x-y)} \;\le\; \tfrac{L_1}{2}\,\Norm{x-y}^{2},
\qquad
\Norm{\nabla F(x) - A}_{op} \le L_1\Norm{x-x^{*}}.
\label{eq:app-aln-l1:R-bounds}
\end{equation}

\paragraph{Итерация VIJI-Restarted на сегменте.}
Согласно~\S\ref{sec:viji:setup} главы~\ref{ch:sp-broyden-viji} на внутренней
итерации сегмента
\begin{align}
x_k     &= v_k - (B_k + \beta_k I)^{-1}\, F(v_k), \label{eq:app-aln-l1:x}\\
v_{k+1} &= (1 - \theta_{k+1})\, x_k + \theta_{k+1}\, v_k,
\quad \theta_{k+1} \in [0,1], \label{eq:app-aln-l1:v}\\
s_k     &= x_k - v_k, \quad y_k = F(x_k) - F(v_k),
\quad B_{k+1} = \mathsf{SP}(B_k, s_k, y_k), \label{eq:app-aln-l1:upd}
\end{align}
с $\beta_k \ge 0$, выбираемым по правилу VIJI
$\beta_k = \max(\beta_{\min}, \eta L_1\Norm{s_{k-1}})$, и шагом следующего
предиктора $d_{k+1} := x_{k+1} - v_{k+1}$.

\paragraph{Обозначения.}
$e_k := v_k - x^{*}$, $\xi_k := x_k - x^{*}$, $\theta := \theta_{k+1}$,
$Z_k := (B_{k+1} + \beta_{k+1} I)^{-1}$, $E_k := B_k - \Jstar$. По
замечанию~\ref{rem:step-ndg-sm} в терминальной фазе равномерно
\begin{equation}
\Norm{Z_k}_{op} = O(1),
\qquad
\Norm{s_k} = O(\Norm{e_k}),
\qquad
\Norm{\xi_k} = O(\Norm{e_k}).
\label{eq:app-aln-l1:basic-bounds}
\end{equation}
Из~\eqref{eq:app-aln-l1:x} следует $F(v_k) = -(B_k + \beta_k I)\,s_k$,
из~\eqref{eq:app-aln-l1:upd}~--- секущее условие $B_{k+1}\,s_k = y_k$.

\paragraph{Ведущая прошлая ошибка.}
Вводим окно глубины $p+1$:
\begin{equation}
\bar e_k \;:=\; \max_{0 \le j \le p+1}\Norm{e_{k-j}}.
\label{eq:app-aln-l1:bar-e}
\end{equation}
В терминальной фазе VIJI-Restarted при Q-линейной сходимости
$\Norm{e_{k+1}} \le \rho_v\Norm{e_k}$, $\rho_v < 1$
\textup{(}\cite[Theorem~3.4]{Agafonov2024VIJI}\textup{)} последовательность
$\{\Norm{e_j}\}$ монотонно убывает, поэтому $\bar e_k = \Norm{e_{k-p-1}}$
и $\Norm{e_{k-j}} \le \bar e_k$ для всех $j \in \{0,\ldots,p+1\}$.

\paragraph{Усреднённые якобианы.}
\begin{equation}
J_k^{(s)} := \int_0^1 \nabla F(v_k + t\, s_k)\,dt,
\qquad
\widetilde J_k := \int_0^1 \nabla F\bigl(v_k + t(1-\theta)\,s_k\bigr)\,dt.
\label{eq:app-aln-l1:avg-J}
\end{equation}
По интегральной формуле $y_k = J_k^{(s)}\,s_k$ и $F(v_{k+1}) - F(v_k) =
(1-\theta)\,\widetilde J_k\,s_k$. Из~\eqref{eq:app-aln-l1:R-bounds}
\begin{equation}
\Norm{(\widetilde J_k - J_k^{(s)})\,s_k} \;=\; O(\Norm{s_k}^{2}),
\qquad
\Norm{J_k^{(s)} - \Jstar}_{op} \;=\; O(\Norm{e_k}).
\label{eq:app-aln-l1:Jbar}
\end{equation}

\section{Тождество разложения шага}
\label{sec:app-aln-l1:identity}

\begin{lemma}[Тождество разложения $d_{k+1}$]
\label{lem:app-aln-l1:identity}
Положим $r_k^{(F)} := \widetilde J_k\,s_k - B_{k+1}\,s_k$. Тогда
\begin{equation}
d_{k+1} \;=\; \theta\,s_k \;+\; (\beta_k - \theta\,\beta_{k+1})\,Z_k\,s_k
\;-\; Z_k\,(B_{k+1} - B_k)\,s_k \;-\; (1-\theta)\,Z_k\,r_k^{(F)},
\label{eq:app-aln-l1:dk-decomp}
\end{equation}
причём $\Norm{r_k^{(F)}} = O(\Norm{s_k}^{2})$.
\end{lemma}

\begin{proof}
Из $F(v_{k+1}) - F(v_k) = (1-\theta)\widetilde J_k\,s_k$ и определения
$r_k^{(F)}$:
$F(v_{k+1}) = F(v_k) + (1-\theta)(B_{k+1}\,s_k + r_k^{(F)})
= -(B_k+\beta_k I)s_k + (1-\theta)B_{k+1}s_k + (1-\theta)r_k^{(F)}$.
Подставляя $-B_k = -B_{k+1} + (B_{k+1}-B_k)$ и добавляя/вычитая
$\theta\beta_{k+1}\,s_k$, получаем
$F(v_{k+1}) = -\theta(B_{k+1}+\beta_{k+1}I)s_k + (\theta\beta_{k+1}-\beta_k)s_k
+ (B_{k+1}-B_k)s_k + (1-\theta)r_k^{(F)}$. Применяя $-Z_k$ и используя
$d_{k+1} = -Z_k F(v_{k+1})$, получаем~\eqref{eq:app-aln-l1:dk-decomp}.
Оценка $r_k^{(F)}$: разлагая $r_k^{(F)} = (\widetilde J_k - J_k^{(s)})s_k +
(J_k^{(s)}s_k - y_k) + (y_k - B_{k+1}s_k)$, второе слагаемое равно нулю
по интегральной формуле, третье~--- по секущему условию, первое
оценивается~\eqref{eq:app-aln-l1:Jbar}.
\end{proof}

При выборе $\alpha_k := \theta_{k+1}$ из~\eqref{eq:app-aln-l1:dk-decomp}
\begin{equation}
d_{k+1} - \alpha_k\,s_k
\;=\; (\beta_k - \theta\beta_{k+1})\,Z_k\,s_k
\;-\; Z_k\,(B_{k+1} - B_k)\,s_k \;-\; (1-\theta)\,Z_k\,r_k^{(F)}.
\label{eq:app-aln-l1:residual}
\end{equation}

\section{Линейная декорреляция секущей: разложение и доказательство}
\label{sec:app-aln-l1:secant-rk}

В этом разделе строится разложение текущей секущей $s_k$ относительно
$p+1$ предшествующих и доказывается ключевая
лемма~\ref{lem:app-aln-l1:secant-rk}~--- линейная декорреляция
ортогональной компоненты $r_k$~--- из стандартного условия
асимптотического выравнивания направлений ошибки в окне
\textup{(}\ref{as:dir-align}\textup{)}. Q-линейная сходимость
$\Norm{e_k}\le\rho_v\Norm{e_{k-1}}$ с $\rho_v < 1$~--- стандартное
свойство VIJI-итерации в терминальной фазе \textup{(}\cite[Theorem~3.4]{Agafonov2024VIJI};
ср.\,\cite[Definition~9.1.5]{dennis1996}\textup{)}~--- используется только
для контроля константы выравнивания и для ограниченности отношений
$\Norm{e_{k-j}}\le\rho_v^{-j}\Norm{e_{k-j-1}}$ при оценке шага
SP-обновления.

\paragraph{Разложение секущей.}
В обозначениях теоремы~\ref{thm:bd} положим
\begin{equation}
S_-^{(k)} := [\,s_{k-1}\mid s_{k-2}\mid\cdots\mid s_{k-p-1}\,] \in\RR^{n\times(p+1)},
\quad
\Pi_k^{<} := S_-^{(k)} \bigl(G_-^{(k)}\bigr)^{-1} \bigl(S_-^{(k)}\bigr)^{\top},
\quad
G_-^{(k)} := (S_-^{(k)})^{\top} S_-^{(k)},
\label{eq:app-aln-l1:Pi-past}
\end{equation}
ортогональный проектор на подпространство $\mathrm{col}\,S_-^{(k)}$. Тогда
\begin{equation}
s_k = \Pi_k^{<}\,s_k + r_k,
\qquad
r_k := (I - \Pi_k^{<})\,s_k \;\perp\; \mathrm{col}\,S_-^{(k)}.
\label{eq:app-aln-l1:sk-decomp}
\end{equation}

\begin{assumption}[Асимптотическое выравнивание направлений ошибки]
\label{as:dir-align}
Существуют константа $\omega_h \ge 0$, зависящая только от
$\mu, L_1, \eta, M_{0}, \rho_v$, и единичный вектор
$\hat u_k \in \RR^{n}$, такие что в терминальной фазе для всех $k$
и всех $j \in \{0,1,\ldots,p+1\}$
\begin{equation}
\bigl\|\hat e_{k-j} - \hat u_k\bigr\| \;\le\; \omega_h\,\bar e_k,
\qquad \hat e_{j} := \frac{e_{j}}{\Norm{e_{j}}}.
\label{eq:app-aln-l1:dir-align}
\end{equation}
\end{assumption}

\paragraph{Природа условия.}
Условие~\eqref{eq:app-aln-l1:dir-align}~--- \emph{структурное
предположение о траектории} VIJI-итераций, того же класса, что
классические направленные условия в анализе квази-Ньютоновских
методов: например, $\Norm{(B_k-\Jstar)s_k}/\Norm{s_k}\to 0$
Дэнниса--Море \textup{(}\cite[Theorem~8.2.4]{dennis1996}\textup{)}
или условие выравнивания секущих в анализе глобальной сходимости
BFGS Пауэлла \textup{(}\cite{powell1971}\textup{)}. В отличие от
условий на оператор $\Jstar, L_1, \mu$ или на схему $\beta_k, \eta$,
оно касается поведения числовой последовательности и не выводимо
непосредственно из~\ref{as:sm}~--- его принятие в качестве условия
теоремы~\ref{thm:app-aln-l1:main} соответствует стандартной
методологии quasi-Newton литературы.

\paragraph{Правдоподобность.}
Линеаризация даёт $e_{k+1} = (I - T_k)e_k + O(\Norm{e_k}^{2})$ с
$T_k := (B_k+\beta_k I)^{-1}\Jstar$, и нормированная динамика
$\hat e_{k+1}$ управляется матрицей $I - T_k$. При простом по модулю
наибольшем собственном значении $\Jstar$ \textup{(}генерическое условие
на оператор\textup{)} степенная итерация обеспечивает выравнивание
$\hat e_k$ вдоль доминирующего собственного направления;
условие~\eqref{eq:app-aln-l1:dir-align}~--- его количественная
формализация в окне глубины $p+1$. Формальный вывод
\eqref{eq:app-aln-l1:dir-align} из спектрального условия на $\Jstar$
при $T_k\to I$ требует анализа slowly-varying matrices в духе Пауэлла
\textup{(}\cite{powell1971}\textup{)} и в общем несимметричном случае
остаётся открытой проблемой~--- зафиксирован
гипотезой~\ref{conj:alignment}. Эмпирически в~\S\ref{sec:viji:numerics}
выравнивание достигается за 2--4 итерации после старта сегмента с
$\omega_h \sim 1$ для всех испытанных сильно монотонных задач.

\begin{lemma}[Медленное изменение $T_k$]
\label{lem:app-aln-l1:slow-T}
В предположениях~\ref{as:sm} и Q-линейной сходимости
$\Norm{e_k}\le\rho_v\Norm{e_{k-1}}$ \textup{(}\cite[Theorem~3.4]{Agafonov2024VIJI}\textup{)}
в терминальной фазе для $j \in \{0,1,\ldots,p+1\}$
\begin{equation}
\Norm{T_{k-j} - T_k}_{op} \;=\; O(\bar e_k).
\label{eq:app-aln-l1:slow-T}
\end{equation}
\end{lemma}

\begin{proof}
Тождество
$T_{k-j} - T_k = (B_{k-j}+\beta_{k-j}I)^{-1}\bigl[(B_k+\beta_k I)
- (B_{k-j}+\beta_{k-j}I)\bigr](B_k+\beta_k I)^{-1}\Jstar$.
Из~\ref{as:sm} обе обратные матрицы равномерно ограничены:
$\Norm{(B_{l}+\beta_{l}I)^{-1}}_{op} = O(1)$. Стандартная оценка для
SP-обновления при адаптивно ограниченном
$\mathrm{cond}(G_-^{(l)})\le\tau$ даёт
$\Norm{B_{l+1}-B_{l}}_{op} = O(\Norm{s_{l}}) = O(\Norm{e_{l}})$
через~\eqref{eq:app-aln-l1:basic-bounds}, поэтому
$\Norm{B_k - B_{k-j}}_{op} \le \sum_{i=1}^{j}O(\Norm{e_{k-i}})
= O((p+1)\bar e_k) = O(\bar e_k)$. Аналогично
$|\beta_k - \beta_{k-j}| = \eta L_1\bigl|\Norm{s_{k-1}}
- \Norm{s_{k-j-1}}\bigr| = O(\bar e_k)$. Совмещая,
получаем~\eqref{eq:app-aln-l1:slow-T}.
\end{proof}

\begin{lemma}[Линейная декорреляция секущей]
\label{lem:app-aln-l1:secant-rk}
В предположениях~\ref{as:sm}, \ref{as:step-ndg}, \ref{as:dir-align} и
Q-линейной сходимости \textup{(}\cite[Theorem~3.4]{Agafonov2024VIJI}\textup{)}
существует константа $\omega_r \ge 0$, зависящая только от
$\mu, L_1, \eta, M_{0}, c_d, \omega_h, \rho_v, p$, такая что в
терминальной фазе для всех $k$
\begin{equation}
\Norm{r_k} \;\le\; \omega_{r}\,\bar e_k\,\Norm{s_k}.
\label{eq:app-aln-l1:secant-rk}
\end{equation}
Эквивалентно: $\Norm{r_k}/\Norm{s_k} = O(\bar e_k)$.
\end{lemma}

\begin{proof}
\emph{Шаг 1 \textup{(}связь $s_{k-j}$ с $e_{k-j}$\textup{)}.}
По $L_1$-гладкости $F(v_{k-j}) = \Jstar e_{k-j} + O(\Norm{e_{k-j}}^{2})$,
поэтому
\begin{equation}
s_{k-j} \;=\; -T_{k-j}\,e_{k-j} + O(\Norm{e_{k-j}}^{2}),
\quad j = 0,1,\ldots,p+1.
\label{eq:app-aln-l1:s-from-e}
\end{equation}

\emph{Шаг 2 \textup{(}подмена $T_{k-j}\to T_k$ и $\hat e_{k-j}\to\hat u_k$\textup{)}.}
По лемме~\ref{lem:app-aln-l1:slow-T} и~\ref{as:dir-align}:
\begin{align*}
T_{k-j}\,e_{k-j}
&= T_k\,e_{k-j} + (T_{k-j}-T_k)\,e_{k-j}
\;=\; \Norm{e_{k-j}}\,T_k\,\hat e_{k-j}
+ O(\bar e_k\,\Norm{e_{k-j}}) \\
&= \Norm{e_{k-j}}\,T_k\,\hat u_k
+ O(\bar e_k\,\Norm{e_{k-j}}).
\end{align*}
Подставляя в~\eqref{eq:app-aln-l1:s-from-e}, получаем общее ведущее
направление $u^{*}_{k} := T_k\,\hat u_k$ для всех секущих окна:
\begin{equation}
s_{k-j} \;=\; -\Norm{e_{k-j}}\,u^{*}_{k} \;+\; w_{k-j},
\qquad
\Norm{w_{k-j}} = O(\bar e_k\,\Norm{e_{k-j}}).
\label{eq:app-aln-l1:s-collinear}
\end{equation}

\emph{Шаг 3 \textup{(}оценка $r_k$\textup{)}.}
При $j=1$ из~\eqref{eq:app-aln-l1:s-collinear}:
$\Norm{e_{k-1}}\,u^{*}_{k} = -s_{k-1} + w_{k-1}$. Умножая на
$\Norm{e_k}/\Norm{e_{k-1}}$ и подставляя $-\Norm{e_k}u^{*}_{k}
= s_k - w_k$ из~\eqref{eq:app-aln-l1:s-collinear} при $j=0$:
\begin{equation*}
s_k \;=\; \frac{\Norm{e_k}}{\Norm{e_{k-1}}}\,s_{k-1}
\;+\; \Bigl(w_k - \frac{\Norm{e_k}}{\Norm{e_{k-1}}}\,w_{k-1}\Bigr).
\end{equation*}
Первое слагаемое лежит в $\mathrm{col}\,S_-^{(k)}$, поэтому проектор
$I-\Pi_k^{<}$ его убивает:
\begin{align*}
\Norm{r_k} = \Norm{(I-\Pi_k^{<})\,s_k}
&\le \Norm{w_k} + \frac{\Norm{e_k}}{\Norm{e_{k-1}}}\Norm{w_{k-1}} \\
&= O(\bar e_k\,\Norm{e_k})
+ \frac{\Norm{e_k}}{\Norm{e_{k-1}}}\cdot O(\bar e_k\,\Norm{e_{k-1}})
\;=\; O(\bar e_k\,\Norm{e_k}).
\end{align*}

\emph{Шаг 4 \textup{(}перевод в форму~\eqref{eq:app-aln-l1:secant-rk}\textup{)}.}
По~\eqref{eq:app-aln-l1:basic-bounds} $\Norm{s_k}\asymp\Norm{e_k}$,
$\Norm{e_k}\le\bar e_k$ тривиально, и потому
$\Norm{r_k} = O(\bar e_k\Norm{e_k}) = O(\bar e_k\Norm{s_k})$.
\end{proof}

\paragraph{Класс задач.}
\ref{as:dir-align} автоматически выполнено
\textup{(}т.~е.~следует из~\ref{as:sm}, \cite[Theorem~3.4]{Agafonov2024VIJI}
плюс условие общего положения на $\Jstar$\textup{)} для следующих классов:
\begin{itemize}\itemsep0pt
  \item Сильно монотонные операторы с $C^{2}$-гладким $\nabla F$
    и простым по модулю наибольшим собственным значением $\Jstar$.
    Кратные доминирующие собственные значения~--- мера-нуль в
    пространстве матриц при фиксированных $\mu,L_1$.
  \item Седловые задачи $\min_x\max_y\Phi(x,y)$ со строгим решением
    второго порядка \textup{(}$\Phi_{xx}\succ 0$, $\Phi_{yy}\prec 0$,
    $\Phi_{xy}$ полного ранга\textup{)}: блочный $\Jstar$ в общем
    положении имеет невырожденный спектр.
  \item Конечно-разностные дискретизации сильно монотонных вариационных
    неравенств в эллиптическом или параболическом режиме~--- спектр
    $\Jstar$ дискретен и без резонансов почти всегда.
\end{itemize}
Граничные случаи \textup{(}изотропные квадратики $\Jstar = \mu I$~---
$r_k=0$ тривиально; $\Jstar$ с нормальной структурой и круговой
симметрией спектра~--- мера-нуль\textup{)} либо не нуждаются
в~\eqref{eq:app-aln-l1:secant-rk}, либо не представляют содержательного интереса.

\paragraph{Замечание о статусе.}
Условие~\ref{as:dir-align}~--- единственное предположение о
траектории, не выводимое из условий на оператор и схему. Введение
спектрального условия на $\Jstar$ \textup{(}простое доминирующее
собственное значение\textup{)} даёт правдоподобность
\ref{as:dir-align}, но не его вывод: переход от спектральной
структуры $\Jstar$ к количественному выравниванию $\hat e_k$ за
окно $p+1$ через степенную итерацию с матрицей $I - T_k$, $T_k\to I$,
требует анализа slowly-varying matrices в духе Пауэлла
\textup{(}\cite{powell1971}\textup{)} и для несимметричных
SP-операторов в общем монотонном случае остаётся открытой
проблемой. Поэтому раздельное введение спектрального условия не
устраняет траекторного характера~\ref{as:dir-align}, а лишь
переносит его в гипотезу о выводе. В настоящем
приложении~\ref{as:dir-align} принято как структурное условие;
полный анализ зафиксирован
гипотезой~\ref{conj:alignment} главы~\ref{ch:sp-broyden-viji} и
предметом отдельной публикации.

\section{Оценка \texorpdfstring{$\Norm{(B_k - \Jstar)\,s_k}$}{|(B_k-J*) s_k|}}
\label{sec:app-aln-l1:dm-bound}

Цель раздела~--- получить из леммы~\ref{lem:app-aln-l1:secant-rk}
строгую оценку $\Norm{E_k s_k} = O(\bar e_k\Norm{s_k})$, нужную в
доказательстве теоремы~\ref{thm:app-aln-l1:main}.

\begin{lemma}[$E_k$ зануляется на прошлых секущих]
\label{lem:app-aln-l1:past-secant}
$B_k$ удовлетворяет $B_k\,s_j = y_j$ для $j = k-1,\ldots,k-p-1$
\textup{(}свойство SP-обновления\textup{)}, поэтому
\begin{equation}
\Norm{E_k\,s_j} \;=\; \Norm{(J_j^{(s)} - \Jstar)\,s_j}
\;=\; O(\Norm{e_j}\,\Norm{s_j}),
\qquad j = k-1,\ldots,k-p-1.
\label{eq:app-aln-l1:past-Esj}
\end{equation}
\end{lemma}

\begin{proof}
$E_k s_j = (B_k - \Jstar)s_j = y_j - \Jstar s_j = (J_j^{(s)} - \Jstar)s_j$
по интегральной формуле. Оценка~$O(\Norm{e_j}\Norm{s_j})$~---
из~\eqref{eq:app-aln-l1:Jbar}.
\end{proof}

\begin{proposition}[Норма $E_k\Pi_k^{<}$]
\label{prop:app-aln-l1:Ek-past}
В предположениях~\ref{as:sm} в терминальной фазе VIJI-Restarted
с адаптивной глубиной SP-Broyden, обеспечивающей
$\mathrm{cond}(G_-^{(k)})\le\tau$,
\begin{equation}
\Fro{E_k\,\Pi_k^{<}} \;=\; O(\bar e_k).
\label{eq:app-aln-l1:Ek-past}
\end{equation}
\end{proposition}

\begin{proof}
Положим $V := E_k\,S_-^{(k)}$. По стандартной оценке:
$\Fro{E_k\,\Pi_k^{<}}^{2} = \tr((G_-^{(k)})^{-1}V^{\top}V) \le
\Fro{V}^{2}/\sigma_{\min}(G_-^{(k)})$. По~\eqref{eq:app-aln-l1:past-Esj}
$\Fro{V}^{2} = O(\sum_{j=1}^{p+1}\Norm{e_{k-j}}^{2}\Norm{s_{k-j}}^{2}) =
O(\bar e_k^{2}\Fro{S_-^{(k)}}^{2})$ \textup{(}из монотонности
$\Norm{e_{k-j}} \le \bar e_k$\textup{)}. Адаптивный порог:
$\sigma_{\min}(G_-^{(k)}) \ge \Fro{S_-^{(k)}}^{2}/((p+1)\tau)$.
Совмещая, $\Fro{E_k\,\Pi_k^{<}}^{2} = O(\bar e_k^{2})$.
\end{proof}

\begin{lemma}[Линейная скорость Дэнниса--Море]
\label{lem:app-aln-l1:dm-rate}
В предположениях~\ref{as:sm}, \ref{as:step-ndg}, \ref{as:dir-align},
теоремы~\ref{thm:bd} с адаптивной глубиной SP-Broyden и Q-линейной
сходимости \textup{(}\cite[Theorem~3.4]{Agafonov2024VIJI}\textup{)}
выполнено
\begin{equation}
\Norm{(B_k - \Jstar)\,s_k} \;=\; O(\bar e_k\,\Norm{s_k}).
\label{eq:app-aln-l1:dm-rate}
\end{equation}
\end{lemma}

\begin{proof}
По~\eqref{eq:app-aln-l1:sk-decomp},
$E_k s_k = E_k\Pi_k^{<}\,s_k + E_k r_k$. Первое слагаемое:
$\Norm{E_k\Pi_k^{<}\,s_k} \le \Fro{E_k\Pi_k^{<}}\Norm{s_k} = O(\bar e_k\Norm{s_k})$
по предложению~\ref{prop:app-aln-l1:Ek-past}. Второе слагаемое: по
лемме~\ref{lem:app-aln-l1:secant-rk} и $\Norm{E_k}_{op}\le M_0$,
$\Norm{E_k r_k} \le M_0\omega_r\bar e_k\Norm{s_k} = O(\bar e_k\Norm{s_k})$.
Складывая, получаем~\eqref{eq:app-aln-l1:dm-rate}.
\end{proof}

\begin{proposition}[Линейный случай: $L_1=0$ влечёт~\eqref{eq:app-aln-l1:dm-rate} безусловно]
\label{prop:app-aln-l1:dm-quad}
При $L_1 = 0$ в~\eqref{eq:app-aln-l1:R-bounds}
$\Norm{(B_k - \Jstar)s_j} = 0$ для всех $j$ из секущей истории
$\{k-1,\ldots,k-p-1\}$, и
\begin{equation}
\sum_{k=k_0}^{\infty} \frac{\Norm{(B_k-\Jstar)\,s_k}^{2}}{\Norm{s_k}^{2}}
\;\le\; \Fro{B_{k_0} - \Jstar}^{\,2}.
\label{eq:app-aln-l1:dm-quad-summable}
\end{equation}
В частности $\Norm{(B_k-\Jstar)s_k}/\Norm{s_k} \to 0$ \emph{без обращения к
\ref{as:dir-align}}; при $p+1 \ge n$ и линейно независимых начальных
секущих сегмента~--- конечно-шаговое равенство $E_k = 0$ для $k \ge k_0+n$.
\end{proposition}

\begin{proof}
При $L_1 = 0$ оператор $F$ линеен, $y_j = A s_j$, и из секущего условия
$E_k s_j = 0$ для $j$ из истории. В обозначениях теоремы~\ref{thm:bd}:
$Y_p - \Jstar S_p = 0$, поэтому $\Delta_k = 0$ в~\eqref{eq:Delta_def}, и
точное разложение~\eqref{eq:exact_decomp} вырождается в монотонное убывание
$\Fro{E_{k+1}}^{2} = \Fro{E_k}^{2} - \Fro{E_k\Pi_k}^{2}$. Телескопирование
$\sum_{k\ge k_0}\Fro{E_k\Pi_k}^{2} \le \Fro{E_{k_0}}^{2}$ и неравенство
$\Norm{E_k s_k}^{2}/\Norm{s_k}^{2} \le \Fro{E_k\Pi_k}^{2}$
дают~\eqref{eq:app-aln-l1:dm-quad-summable}. Конечно-шаговое равенство
при $p+1\ge n$~--- из $\Pi_k = I$ и $\Fro{E_{k+1}}^{2} = 0$.
\end{proof}

\section{Заключительная теорема}
\label{sec:app-aln-l1:close}

\begin{theorem}[Лемма~\ref{lem:alignment} в общем $L_1$-гладком случае]
\label{thm:app-aln-l1:main}
В предположениях~\ref{as:sm}, \ref{as:step-ndg}, \ref{as:dir-align},
Q-линейной сходимости \textup{(}\cite[Theorem~3.4]{Agafonov2024VIJI}\textup{)}
и при стандартном правиле VIJI ($\theta_{k+1}\in[0,1]$) существуют
константы $C_a, C_b > 0$, зависящие только от
$\mu, L_1, \eta, M_0, c_d, \omega_h, \rho_v, p$, такие что для всех $k$
в терминальной фазе \textup{(}$\bar e_k \le
\varepsilon_{\mathrm{ndg}}$\textup{)} с $\alpha_k := \theta_{k+1}$
\begin{equation}
\Norm{d_{k+1} - \alpha_k\,s_k}
\;\le\; C_a\,\bar e_k\Norm{d_{k+1}}
\;+\; C_b\,\Norm{d_{k+1}}^{2}.
\label{eq:app-aln-l1:final}
\end{equation}
В правой части $\bar e_k = \max_{0\le j\le p+1}\Norm{v_{k-j}-x^{*}}$~---
ведущая прошлая ошибка по окну SP-Broyden глубины $p+1$
\textup{(}связь с условием~\eqref{eq:cond14} см.~ниже,
п.~<<Связь с условием~\eqref{eq:cond14}>>\textup{)}.
\end{theorem}

\begin{proof}
Оценим в~\eqref{eq:app-aln-l1:residual} каждое из трёх слагаемых;
все $O(1)$-константы зависят от $\mu,L_1,\eta,M_0,c_d,\omega_r$.

\emph{Регуляризатор $\beta$.} Правило VIJI: $\beta_k = \eta L_1\Norm{s_{k-1}}$,
$\beta_{k+1} = \eta L_1\Norm{s_k}$, оба $\ge 0$, поэтому
$|\beta_k - \theta\beta_{k+1}| \le \eta L_1\max(\Norm{s_{k-1}},\Norm{s_k})$.
По~\eqref{eq:app-aln-l1:basic-bounds} и монотонности $\{\Norm{e_j}\}$:
$|\beta_k-\theta\beta_{k+1}| = O(\Norm{e_{k-1}})$, $\Norm{Z_k s_k} =
O(\Norm{e_k})$. Применением~\ref{as:step-ndg} ($\Norm{e_k}\le c_d^{-1}\Norm{d_{k+1}}$)
и $\Norm{e_{k-1}}\le\bar e_k$:
\begin{equation}
\Norm{(\beta_k-\theta\beta_{k+1})Z_k s_k} \;=\; O(\bar e_k\Norm{d_{k+1}}).
\label{eq:app-aln-l1:residual-beta}
\end{equation}

\emph{Секущая разность $B_{k+1}-B_k$.} Из секущего условия и интегральной
формулы $(B_{k+1}-B_k)s_k = y_k - B_k\,s_k = -E_k\,s_k + (J_k^{(s)} - \Jstar)s_k$.
По лемме~\ref{lem:app-aln-l1:dm-rate} первое слагаемое~$O(\bar e_k\Norm{s_k})$,
по~\eqref{eq:app-aln-l1:Jbar} второе~$O(\Norm{e_k}\Norm{s_k}) =
O(\bar e_k\Norm{s_k})$. Через~\ref{as:step-ndg}
\begin{equation}
\Norm{Z_k(B_{k+1}-B_k)s_k} \;=\; O(\bar e_k\Norm{d_{k+1}}).
\label{eq:app-aln-l1:residual-B}
\end{equation}

\emph{Якобиан $r_k^{(F)}$.} По лемме~\ref{lem:app-aln-l1:identity}
$\Norm{Z_k r_k^{(F)}} = O(\Norm{s_k}^{2}) = O(\Norm{e_k}^{2})$, через
\ref{as:step-ndg}
\begin{equation}
\Norm{(1-\theta)Z_k r_k^{(F)}} \;=\; O(\Norm{d_{k+1}}^{2}).
\label{eq:app-aln-l1:residual-F}
\end{equation}

Складывая~\eqref{eq:app-aln-l1:residual-beta}--\eqref{eq:app-aln-l1:residual-F}
и применяя неравенство треугольника к~\eqref{eq:app-aln-l1:residual},
получаем~\eqref{eq:app-aln-l1:final}.
\end{proof}

\section{Замечания}
\label{sec:app-aln-l1:remarks}

\paragraph{Связь с условием~\eqref{eq:cond14}.}
Оценка~\eqref{eq:app-aln-l1:final} отличается от исходной формы
леммы~\ref{lem:alignment} тем, что в первом слагаемом стоит \emph{прошлая}
ошибка $\bar e_k = \Norm{v_{k-p-1}-x^{*}}$ вместо текущей
$\Norm{v_k-x^{*}}$. Применение в~\S\ref{sec:viji:theorem}~--- переход от
$\Norm{(B_{k+1}-J_k)s_k}$ к левой части~\eqref{eq:cond14}~--- проходит
после одного дополнительного шага: связки $\bar e_k = O(\Norm{d_{k+1}})$.
Достаточные условия:
\begin{enumerate}\itemsep0pt
  \item \emph{Двусторонняя Q-линейная сходимость} в терминальной фазе:
    $\underline\rho_v\Norm{e_{k-1}}\le\Norm{e_k}\le\rho_v\Norm{e_{k-1}}$
    с $\underline\rho_v > 0$, откуда $\bar e_k = \Norm{e_{k-p-1}}
    \le \underline\rho_v^{-(p+1)}\Norm{e_k} = O(\Norm{e_k}) =
    O(\Norm{d_{k+1}})$ через~\ref{as:step-ndg}. Это
    \emph{дополнительное} условие на динамику, более сильное, чем
    стандартная Q-линейная оценка~\cite[Theorem~3.4]{Agafonov2024VIJI},
    и в работе как ассумпция не вводится.
  \item \emph{Окно-локальная невырожденность шага}: применение
    \ref{as:step-ndg} ко всем недавним итерациям окна плюс ограниченность
    отношений $\Norm{d_{k-q+1}}/\Norm{d_{k+1}}$~--- эквивалентно
    предыдущему пункту.
\end{enumerate}
Без любого из этих допущений лемма~\ref{lem:alignment} даёт оценку с
$\bar e_k$ вместо $\Norm{v_k-x^{*}}$, и закрытие~\eqref{eq:cond14}
требует лишь $\bar e_k\to 0$ \textup{(}обеспечено Q-линейной сходимостью\textup{)};
переход к квадратичной по $\Norm{d_{k+1}}$ форме требует двусторонней
оценки.
